% English translation of chapter2.tex lines 1133--1244.
% Source: Methods of Algebra, Volume 2, commit
% 9a5803ff77dd3257484cb177f851a73770a59dd3 (CC BY 4.0).
% stable-unit-id: o014.aljabr2.chapter2.simplicity-and-semisimplicity

% segment-id: o014.li2.chapter2.simplicity-and-semisimplicity.h001
\section{Simple and Semisimple Objects}\label{sec:semisimple}
\hypertarget{o014.aljabr2.chapter2.simplicity-and-semisimplicity}{ }

% segment-id: o014.li2.chapter2.simplicity-and-semisimplicity.p001
Throughout this section, let $\mathcal{A}$ be an abelian category.

% segment-id: o014.li2.chapter2.simplicity-and-semisimplicity.c001
\begin{convention}\label{con:direct-sum}
	\index{direct sum}
	By convention, the coproduct $\coprod_{i \in I} X_i$ of a family of
	objects $(X_i)_{i \in I}$ in an abelian category, whenever it exists, is
	written $\bigoplus_{i \in I} X_i$ and called their \emph{direct sum}. If
	$I$ is finite, all the conventions of Definition~\ref{def:additive-category}
	apply.
\end{convention}

% segment-id: o014.li2.chapter2.simplicity-and-semisimplicity.d001
\begin{definition}\index{simple object}
	A nonzero object $X$ of $\mathcal{A}$ is called \emph{simple} if
	$\mathrm{Sub}_X=\left\{0,X\right\}$.
\end{definition}

% segment-id: o014.li2.chapter2.simplicity-and-semisimplicity.p002
Simplicity of $X$ means that for every short exact sequence
$0\to X'\to X\to X''\to0$, exactly one of $X'=0$ and $X'\rightiso X$
holds; equivalently, exactly one of $X\rightiso X''$ and $X''=0$ holds.
Consequently, $X$ is simple in $\mathcal{A}$ if and only if it is simple in
$\mathcal{A}^{\opp}$.

% segment-id: o014.li2.chapter2.simplicity-and-semisimplicity.p003
Notice that $X\neq0$ if and only if $\identity_X\neq0$, and this is also
equivalent to $\End(X)$ being a nonzero ring.

% segment-id: o014.li2.chapter2.simplicity-and-semisimplicity.lm001
\begin{lemma}[Schur's lemma]
	Let $X,Y$ be objects. If $X$ (respectively, $Y$) is simple, then every
	nonzero morphism in $\Hom(X,Y)$ is a monomorphism (respectively, an
	epimorphism). Consequently, if $X$ is simple, then $\End(X)$ is a division
	ring.
\end{lemma}
% segment-id: o014.li2.chapter2.simplicity-and-semisimplicity.pf001
\begin{proof}
	Suppose that $X$ is simple and that $f:X\to Y$ is nonzero. Then necessarily
	$\Ker(f)=0$. The case in which $Y$ is simple follows by duality. A morphism
	that is both a monomorphism and an epimorphism is an isomorphism
	(Proposition~\ref{prop:strict-isomorphism}); taking $X=Y$ therefore shows
	that $\End(X)$ is a division ring.
\end{proof}

% segment-id: o014.li2.chapter2.simplicity-and-semisimplicity.p004
We know that $\mathrm{Sub}_X$ is a bounded modular lattice
(Theorem~\ref{prop:subobject-modularity}). Lemma~\ref{prop:lattice-composition-series}
shows that $X$ has finite length if and only if the partially ordered set
$[0,X]$ has a composition series; such a series is also called a composition
series of $X$. For an object $X$ of finite length, its \emph{length} is
defined by \index{length}
\index[sym1]{lX@$\ell(X)$}
% segment-id: o014.li2.chapter2.simplicity-and-semisimplicity.q001
\[ \ell(X) := \text{the length of }\mathrm{Sub}_X\;\in\Z_{\geq0}; \]
see Definition~\ref{def:lattice-length}. Notice that
$\ell(X)=0\iff X=0$. By convention, $Y\subsetneq Z$ means $Y\subset Z$ and
$Y\neq Z$.

% segment-id: o014.li2.chapter2.simplicity-and-semisimplicity.dt001
\begin{definition-theorem}[Jordan--H\"older theorem]\label{def:JH}
	\index{Jordan--Holder theorem@Jordan--H\"older theorem}
	\index{composition factor}
	\index[sym1]{JH@$\mathrm{JH}(X)$}
	\sloppy
	Let $X$ be an object of finite length, and choose a composition series
	$X=X_0\supsetneq\cdots\supsetneq X_r=0$. The objects $X_i/X_{i+1}$,
	considered up to isomorphism, are called the \emph{composition factors} of
	$X$. They are simple. Their multiset, with multiplicities, is denoted
	$\mathrm{JH}(X)$ and is independent of the chosen composition series.
	\fussy
\end{definition-theorem}
% segment-id: o014.li2.chapter2.simplicity-and-semisimplicity.pf002
\begin{proof}
	The composition factors must be simple, for otherwise
	$X_0\supsetneq\cdots\supsetneq X_r$ would admit a proper refinement. Let
	$X=X'_0\supsetneq\cdots\supsetneq X'_s$ be another composition series.
	Theorem~\ref{prop:lattice-JH} says that it is equivalent to
	$X_0\supsetneq\cdots\supsetneq X_r$. In particular, there is a bijection
	$i\leftrightarrow j$ between their index sets such that the intervals
	$\left[X_{i+1},X_i\right]$ and $\left[X'_{j+1},X'_j\right]$ can be linked by standard
	isomorphisms $[a\wedge b,a]\rightiso[b,a\vee b]$ in a modular lattice;
	see Proposition~\ref{prop:lattice-diamond-isom}.

	So far everything has been phrased in lattice-theoretic language. However,
	Theorem~\ref{prop:Abel-cat-isom-thm}(iii) shows that such interval
	isomorphisms correspond to isomorphisms of quotient objects in
	$\mathcal{A}$,
% segment-id: o014.li2.chapter2.simplicity-and-semisimplicity.q002
	\[ X_i/X_{i+1} \simeq X'_j/X'_{j+1}. \]
	As $i\leftrightarrow j$ varies, this proves that the composition factors
	are independent of the chosen composition series, up to isomorphism and
	reordering.
\end{proof}

% segment-id: o014.li2.chapter2.simplicity-and-semisimplicity.p005
The number of elements of $\mathrm{JH}(X)$, counted with multiplicity, is
exactly $\ell(X)$. Multisets also have a union operation, which adds
multiplicities; it will again be denoted by $\cup$.

% segment-id: o014.li2.chapter2.simplicity-and-semisimplicity.lm002
\begin{lemma}
	Given a short exact sequence $0\to X'\to X\to X''\to0$, the object $X$
	has finite length if and only if both $X'$ and $X''$ do. In that case,
	$\mathrm{JH}(X)=\mathrm{JH}(X')\cup\mathrm{JH}(X'')$, and hence
	$\ell(X)=\ell(X')+\ell(X'')$.
\end{lemma}
% segment-id: o014.li2.chapter2.simplicity-and-semisimplicity.pf003
\begin{proof}
	Theorem~\ref{prop:Abel-cat-isom-thm}(ii) identifies
	$\mathrm{Sub}_{X'}$ and $\mathrm{Sub}_{X''}$ with the intervals $[0,X']$
	and $[X',X]$ in $\mathrm{Sub}_X$, respectively. Thus finite length of $X$
	implies finite length of $X'$ and $X''$. Conversely, suppose that $X'$ and
	$X''$ have finite length, and choose composition series
% segment-id: o014.li2.chapter2.simplicity-and-semisimplicity.q003
	\[ X'=X'_0\supsetneq\cdots\supsetneq X'_r=0,\quad
	X''=X''_0\supsetneq\cdots\supsetneq X''_s=0. \]
	By Theorem~\ref{prop:Abel-cat-isom-thm}(ii), lift every $X''_i$ to a
	subobject $Y_i$ of $X$ satisfying $Y_i\supset X'$. In particular, $Y_0=X$ and
	$Y_s=X'$, so
% segment-id: o014.li2.chapter2.simplicity-and-semisimplicity.q004
	\[ X=Y_0\supsetneq\cdots\supsetneq Y_s\supsetneq X'_1\supsetneq\cdots
	\supsetneq X'_r=0 \]
	is a composition series whose subquotients form
	$\mathrm{JH}(X')\cup\mathrm{JH}(X'')$.
\end{proof}

% segment-id: o014.li2.chapter2.simplicity-and-semisimplicity.d002
\begin{definition}\label{def:semisimple}
	\index{object!semisimple}
	\index{object!split}
	\index{abelian category!semisimple}
	\index{abelian category!split}
	An object $X$ is called
	\begin{itemize}
		% segment-id: o014.li2.chapter2.simplicity-and-semisimplicity.i001
		\item \emph{semisimple} if
		$X=\bigoplus_{i\in I}X_i$ for some family of simple objects
		$(X_i)_{i\in I}$;
		% segment-id: o014.li2.chapter2.simplicity-and-semisimplicity.i002
		\item \emph{split} if every short exact sequence
		$0\to X'\to X\to X''\to0$ splits.
	\end{itemize}
	If every object of $\mathcal{A}$ is semisimple (respectively, split), then
	$\mathcal{A}$ is called a semisimple (respectively, split) abelian
	category.\footnote{Terminology in the literature is not uniform; some
	authors call what is termed a split abelian category here a semisimple
	abelian category.}
\end{definition}

% segment-id: o014.li2.chapter2.simplicity-and-semisimplicity.p006
Many sources require $I$ to be finite in the definition of a semisimple
object.

% segment-id: o014.li2.chapter2.simplicity-and-semisimplicity.lm003
\begin{lemma}
	If $X=\bigoplus_{i=1}^nX_i$, where $X_1,\ldots,X_n$ are simple, then $X$
	has finite length,
	$\mathrm{JH}(X)=\left\{X_1,\ldots,X_n\right\}$, counted with multiplicity, and
	$\ell(X)=n$.
\end{lemma}
% segment-id: o014.li2.chapter2.simplicity-and-semisimplicity.pf004
\begin{proof}
	Consider the composition series
	$\bigoplus_{i=1}^nX_i\supsetneq\bigoplus_{i=1}^{n-1}X_i\supsetneq\cdots
	\supsetneq0$.
\end{proof}

% segment-id: o014.li2.chapter2.simplicity-and-semisimplicity.rem001
\begin{remark}
	For a general semisimple object $X=\bigoplus_{i\in I}X_i$, the
	Noetherian, Artinian, and finite-length properties of $X$ are all
	equivalent to finiteness of $I$. Indeed, by adding (or removing) direct
	summands indexed by $I$, one can readily construct a strictly ascending
	(or strictly descending) chain in $\mathrm{Sub}_X$.
\end{remark}

% segment-id: o014.li2.chapter2.simplicity-and-semisimplicity.p007
We wish to understand the relation between split and semisimple objects. The
argument is similar to the module case \cite[Proposition 6.11.4]{Li1}.

% segment-id: o014.li2.chapter2.simplicity-and-semisimplicity.pr001
\begin{proposition}
	Let $X$ be a split object. Then every subobject and every quotient object of
	$X$ is split. If, in addition, $X$ is Artinian, then $X$ is a semisimple
	object of finite length.
\end{proposition}
% segment-id: o014.li2.chapter2.simplicity-and-semisimplicity.pf005
\begin{proof}
	Given a short exact sequence $0\to X'\to X\to X''\to0$, the hypothesis
	implies $X\simeq X'\oplus X''$, so $X''$ embeds as a subobject of $X$.
	It is therefore enough to show that every subobject $X'$ of $X$ is split.
	Given $X'_0\subset X'$, there is a $Y\subset X$ such that
	$X=X'_0\oplus Y$. It remains to prove
% segment-id: o014.li2.chapter2.simplicity-and-semisimplicity.q005
	\[ X'=X'_0\oplus(Y\cap X'). \]
	This follows from Proposition~\ref{prop:biproduct-sum-intersection}. On the
	one hand, $X'_0\cap(Y\cap X')\subset X'_0\cap Y=0$. On the other hand,
	$\mathrm{Sub}_X$ is modular, and hence
	$X'_0+(Y\cap X')=X'\cap(Y+X'_0)=X'\cap X=X'$. This proves the displayed
	decomposition.

	Now suppose also that $X$ is Artinian. If $X\neq0$, it has a minimal
	nonzero subobject $X_1$, which must be simple, and there is a direct-sum
	decomposition $X=X_1\oplus Y_1$. The object $Y_1$ is still Artinian and
	split. If $Y_1\neq0$, continue by writing $Y_1=X_2\oplus Y_2$, and so on.
	The Artinian condition ensures that the strictly descending chain
	$X\supsetneq Y_1\supsetneq Y_2\cdots$ terminates after finitely
	many steps, giving the required decomposition
	$X=X_1\oplus\cdots\oplus X_n$.
\end{proof}

% segment-id: o014.li2.chapter2.simplicity-and-semisimplicity.p008
We may ask conversely whether semisimple objects are split. For a general
abelian category, a finiteness assumption is again needed.

% segment-id: o014.li2.chapter2.simplicity-and-semisimplicity.pr002
\begin{proposition}\label{prop:ss-decomp}
	For an object $X$, consider the following properties.
	\begin{enumerate}[(i)]
		% segment-id: o014.li2.chapter2.simplicity-and-semisimplicity.i003
		\item $X=\sum_{Y\in\mathcal{F}}Y$, where $\mathcal{F}$ is a finite
		subset of $\mathrm{Sub}_X$ and every $Y\in\mathcal{F}$ is simple;
		% segment-id: o014.li2.chapter2.simplicity-and-semisimplicity.i004
		\item $X=\bigoplus_{Y\in\mathcal{F}}Y$, where $\mathcal{F}$ is a
		finite subset of $\mathrm{Sub}_X$ and every $Y\in\mathcal{F}$ is
		simple;
		% segment-id: o014.li2.chapter2.simplicity-and-semisimplicity.i005
		\item $X$ is split.
	\end{enumerate}
	Then (i) $\implies$ (ii) $\implies$ (iii).
\end{proposition}
% segment-id: o014.li2.chapter2.simplicity-and-semisimplicity.pf006
\begin{proof}
	Repeat the proof of (i) $\implies$ (ii) $\implies$ (iii) in
	\cite[Proposition 6.11.4]{Li1}.
\end{proof}

% segment-id: o014.li2.chapter2.simplicity-and-semisimplicity.rem002
\begin{remark}\label{rem:Grothendieck-ss-decomp}
	To extend Proposition~\ref{prop:ss-decomp} to an infinite subset
	$\mathcal{F}\subset\mathrm{Sub}_X$, one must require $\mathcal{A}$ to be
	a Grothendieck category, as introduced in
	\S~\sourcecrossref{sec:Grothendieck-cat}{2.10}. Every semisimple object in
	a Grothendieck category is then split, and every semisimple Grothendieck
	category is automatically split. Once the relevant definitions are in
	place, the argument is similar to the module case and is left as an
	exercise in this chapter.
\end{remark}
